\providecommand{\Needspace}[1]{}
\Needspace{8\baselineskip}
\item \textbf{Ajuste cuadrático y estimación de un umbral}.
\nopagebreak[4]

\begin{tcolorbox}[colback=blue!2!white,colframe=blue!55!black,title={Enunciado},
                  before skip=3pt,after skip=5pt,boxsep=3pt,top=3pt,bottom=3pt,
                  breakable]
\small
\raggedright
En una calibración de laboratorio se mide la respuesta \(R\) de un detector para
distintas dosis \(d\). Los datos medidos son:
\[
d = [0,\ 1,\ 2,\ 3,\ 4,\ 5,\ 6],
\]
\[
R = [1.2,\ 2.0,\ 4.1,\ 6.9,\ 10.8,\ 16.2,\ 22.1].
\]

Se desea construir un modelo cuadrático
\[
R(d)=ad^2+bd+c
\]
y usarlo para estimar cuándo la respuesta alcanza un umbral de seguridad
\[
R=12.
\]

Realice las siguientes operaciones:

\begin{itemize}
\setlength{\topsep}{4pt}
\setlength{\itemsep}{3pt}
\setlength{\parsep}{0pt}
\item[(a)] Importe \texttt{numpy} y \texttt{matplotlib.pyplot}. Defina los
arreglos \texttt{d} y \texttt{R} usando \texttt{np.array}.
\item[(b)] Use \texttt{np.polyfit(d, R, deg=2)} para obtener los coeficientes
\(a\), \(b\) y \(c\).
\item[(c)] Calcule los valores ajustados sobre los mismos datos medidos y
guárdelos en \texttt{R\_fit}.
\item[(d)] Calcule los residuos y el coeficiente \(R^2\):
\[
\texttt{residuos = R - R\_fit},
\]
\[
SS_{\mathrm{res}} = \sum_i \left(R_i - R_{\mathrm{fit},i}\right)^2,
\qquad
SS_{\mathrm{tot}} = \sum_i \left(R_i - \overline{R}\right)^2,
\]
\[
R^2 = 1 - \frac{SS_{\mathrm{res}}}{SS_{\mathrm{tot}}}.
\]
\item[(e)] Use \texttt{np.roots} para resolver cuando el modelo ajustado alcanza
\(R=12\). Para ello debe resolver
\[
ad^2+bd+c-12=0.
\]
\item[(f)] Filtre las raíces: conserve solo soluciones reales dentro del rango
medido, es decir, entre \(0\) y \(6\).
\item[(g)] Grafique los datos medidos con \texttt{plt.scatter}, el ajuste
cuadrático con \texttt{plt.plot} usando una grilla fina para la dosis, y la
línea horizontal \(R=12\). Para esta grilla fina, cree un arreglo con
\texttt{np.linspace} para la variable independiente. Agregue etiquetas, leyenda
y active la grilla del gráfico.
\item[(h)] Imprima \(a\), \(b\), \(c\), \(R^2\) y la dosis estimada para alcanzar
el umbral. Agregue una frase que explique por qué no basta con tomar cualquier
raíz entregada por \texttt{np.roots}.
\end{itemize}
\end{tcolorbox}

\begin{solution}
\begin{pythonjupyter}
import numpy as np
import matplotlib.pyplot as plt

d = np.array([0, 1, 2, 3, 4, 5, 6])
R = np.array([1.2, 2.0, 4.1, 6.9, 10.8, 16.2, 22.1])

a, b, c = np.polyfit(d, R, deg=2)

R_fit = a*d**2 + b*d + c

residuos = R - R_fit
SS_res = np.sum(residuos**2)
SS_tot = np.sum((R - np.mean(R))**2)
R2 = 1 - SS_res/SS_tot

coef_umbral = [a, b, c - 12]
raices = np.roots(coef_umbral)

soluciones = []
for r in raices:
    if abs(r.imag) < 1e-10 and 0 <= r.real <= 6:
        soluciones.append(r.real)

d_fino = np.linspace(0, 6, 200)
R_fino = a*d_fino**2 + b*d_fino + c

plt.figure(figsize=(7, 5))
plt.scatter(d, R, label="datos medidos")
plt.plot(d_fino, R_fino, "r-", label="ajuste cuadrático")
plt.axhline(12, color="gray", linestyle="--", label="umbral R=12")
plt.xlabel("Dosis d")
plt.ylabel("Respuesta R")
plt.title("Ajuste cuadrático y umbral de seguridad")
plt.legend()
plt.grid(True, alpha=0.3)
plt.show()

print(f"a = {a:.4f}")
print(f"b = {b:.4f}")
print(f"c = {c:.4f}")
print(f"R2 = {R2:.4f}")

if len(soluciones) > 0:
    print(f"Dosis estimada para R=12: {soluciones[0]:.3f}")
else:
    print("No hay una solución real dentro del rango medido")

print("No basta con tomar cualquier raíz de np.roots:")
print("hay que descartar raíces complejas y raíces fuera del rango físico.")
\end{pythonjupyter}
\end{solution}
